% Generated draft: requires mathematical and source review.
\documentclass[11pt]{article}
\usepackage[margin=25mm]{geometry}
\usepackage{fontspec}
\setmainfont{Latin Modern Roman}
\usepackage{amsmath,amssymb,mathtools,unicode-math}
\setmathfont{Latin Modern Math}
\usepackage{enumitem,needspace,xurl,hyperref}
\setlength{\parindent}{0pt}
\setlength{\parskip}{5pt}
\setlength{\emergencystretch}{3em}
\newcommand{\N}{\mathbb{N}}
\newcommand{\Z}{\mathbb{Z}}
\newcommand{\Q}{\mathbb{Q}}
\newcommand{\R}{\mathbb{R}}
\newcommand{\C}{\mathbb{C}}
\newcommand{\F}{\mathbb{F}}
\newcommand{\A}{\mathbb{A}}
\newcommand{\PP}{\mathbb{P}}
\newcommand{\E}{\mathbb{E}}
\newcommand{\eps}{\varepsilon}
\newcommand{\dd}{\,\mathrm{d}}
\begin{document}
\section*{UnsolvedMath ID: 30002876}
\textbf{Source problem code:} catalogue entry
\subsection{Definitions and mathematical statement}
The supplied record does not specify the algebraic group \(G\), its base field or characteristic, or the allowed values of \(r\). Thus its intended scope cannot be recovered completely. The following gives an explicit conditional interpretation in a conventional setting; these hypotheses are not supplied source facts.

Suppose that \(k\) is an algebraically closed field of characteristic \(p>0\), that \(G\) is a connected, simply connected semisimple algebraic group over \(k\), and that \(r\geq 1\) is an integer. Fix a maximal torus \(T\subseteq G\) and a system of positive roots. Let \(X(T)=\operatorname{Hom}(T,\mathbb G_m)\) be the character lattice, let \(X^+\subseteq X(T)\) consist of the weights whose pairing with every simple coroot is nonnegative, and let \(\rho\) be half the sum of the positive roots. For \(\lambda\in X^+\), let \(L(\lambda)\) denote the finite-dimensional simple rational \(G\)-module of highest weight \(\lambda\). In this setting, define the \(r\)-th Steinberg module by
\[
\mathrm{St}_r=L\bigl((p^r-1)\rho\bigr).
\]
For \(\mu\in X^+\), let \(\nabla(\mu)\) denote the dual Weyl, or costandard, rational \(G\)-module of highest weight \(\mu\). A finite-dimensional rational \(G\)-module \(M\) has a good filtration if it admits a finite filtration by rational \(G\)-submodules
\[
0=M_0\subset M_1\subset\cdots\subset M_n=M
\]
such that, for every \(1\leq i\leq n\), there is a weight \(\mu_i\in X^+\) with \(M_i/M_{i-1}\cong\nabla(\mu_i)\).

For fixed parameters in this conditional setting, the source question is to determine exactly the set
\[
\bigl\{\lambda\in X^+ : \mathrm{St}_r\otimes_k L(\lambda)\text{ has a good filtration as a rational }G\text{-module}\bigr\},
\]
where the tensor product has the diagonal \(G\)-action. The intended hypotheses on \(G\), \(k\), \(p\), and \(r\), and whether the classification is requested uniformly over any of these parameters, remain unresolved in the supplied record.
\subsection{Short English statement}
Determine which dominant weights \(\lambda\) make the tensor product of the \(r\)-th Steinberg module with the simple module \(L(\lambda)\) admit a filtration with dual Weyl module factors. The source omits the ambient hypotheses, so the displayed formulation is conditional and requires review against the original report.
\subsection{Sources}
\begin{itemize}
\item Supplied local UnsolvedMath registry record, permanent ID \(30002876\), title ``Steinberg Tensor Products with Good Filtrations,'' Oberwolfach Reports Open Problems collection, exact source problem code \texttt{catalogue entry}. No individual proposer is supplied; the proposed year is recorded as 2015.
\item Supplied report citation: \emph{Cohomology of Finite Groups: Interactions and Applications} (2016), Oberwolfach Reports, pp.\ 1317--1384, DOI \texttt{10.4171/owr/2015/24}. \url{https://doi.org/10.4171/owr/2015/24}. The record describes Oberwolfach Reports as documenting mathematical workshops and their open-problem discussions. The report citation has not been independently consulted here.
\item Supplied literature reference: \emph{On tensoring with the Steinberg representation}. \url{https://arxiv.org/abs/1804.00613}.
\item Supplied literature reference: \emph{Counterexamples to the Tilting and (p,r)-Filtration Conjectures}. \url{https://arxiv.org/abs/1901.06687}.
\item Literature assessment dated exactly \texttt{2026-08-21}: the supplied curation labels the problem open and states that the Steinberg-tensor good-filtration question is settled only in broad special ranges and cases, while the general classification of \(\lambda\) remains open. This is supplied dated curation derived from the OWR record, not an independently verified literature assessment. The record explicitly requests independent verification of time-sensitive status and cited claims.
\item The registry labels the statement \texttt{corrected\_verified} and reports that its mathematical content was preserved while notation, typography, or omitted definitions were repaired for publication. This is a supplied editorial assertion, not mathematical verification of this draft.
\end{itemize}
\end{document}
